\chapter{Costruzione e Risoluzione delle Query Architetturali}
\label{ch:name_resolution}

La fase di \textit{Reference Resolution} (Risoluzione dei
Riferimenti) è il cuore logico dell'analizzatore architetturale. Il
suo scopo è trasformare riferimenti astratti ed esposti ad ambiguità
testuali (i nomi grezzi estratti dall'albero sintattico) in
identificatori univoci e globalmente validi, creando i collegamenti
necessari per il grafo delle dipendenze.

Nelle prime iterazioni del sistema, la risoluzione operava in un
singolo passaggio sincrono (eseguendo il \textit{lookup} non appena
il nodo veniva incontrato). Tale approccio si è rivelato inadeguato
per due motivi fondamentali: l'impossibilità di risolvere i
\textit{forward references} (tipi dichiarati testualmente dopo il
loro utilizzo) e l'incapacità di analizzare catene di chiamate di
metodi (es. \texttt{a.f().g()}).
Per allinearci alla teoria dei compilatori, abbiamo riprogettato la
risoluzione introducendo un \textbf{Query Engine Bifase}, basato sul
principio della \textit{lazy evaluation} (valutazione pigra).

\section{Fase 2a: Costruzione delle Query (Substitution Phase)}

Nella prima sotto-fase ($R_{build}$), l'obiettivo non è scoprire
\textit{dove} si trovi un componente, ma formalizzare
\textit{l'intento logico} della ricerca.

Le variabili locali (es. \texttt{let v: a.b.T}) sono essenziali per
il tracciamento del flusso di controllo, ma non costituiscono nodi
architetturali nel grafo finale. Vengono utilizzate esclusivamente in
questa fase come ponte logico. L'analizzatore attraversa l'IR
avvalendosi di un ambiente lessicale effimero, il \textbf{Symbol
Stack} ($\rho$).
Quando viene dichiarata una variabile locale, l'ambiente $\rho$ ne
registra l'associazione testuale. Quando la variabile viene
utilizzata, il sistema consuma il riferimento sostituendo
sintatticamente il nome locale con la formula matematica
corrispondente al suo tipo, generando un costrutto intermedio
chiamato \texttt{ResolutionQuery($\mathcal{Q}$)}.

\subsection{L'Algebra delle Query ($\mathcal{Q}$)}

Per garantire la massima astrazione, ogni operazione di accesso o
invocazione in qualsiasi linguaggio di programmazione è stata ridotta
alla combinazione di sole tre primitive matematiche. Queste primitive
catturano le direzioni di navigazione (salita e discesa) nel grafo
gerarchico dei moduli.

\begin{itemize}
  \item \texttt{Find(name)} (\textit{Ricerca Ascendente}): Definisce
    il \textbf{punto di partenza assoluto} dell'espressione. È
    necessario perché un'espressione come \texttt{T.x} non ha
    significato finché non si stabilisce a quale \texttt{T} si faccia
    riferimento nell'albero dei moduli. Il \texttt{Find} ordina al
    motore di risoluzione di agganciarsi al contesto lessicale
    corrente e "risalire" progressivamente la gerarchia degli scope
    (Lexical Scope Climbing) finché non individua la prima entità chiamata
    \texttt{name}.
  \item \texttt{Extract(query, name)} (\textit{Discesa Strutturale}):
    Rappresenta l'operatore di accesso ai membri (es. il punto
    \texttt{.} o i due punti \texttt{::}). L'operatore riceve in
    input una sub-query (il contesto padre) e il nome del membro
    figlio da cercare. Valuta la sub-query per ottenere l'indirizzo
    del genitore, dopodiché discende di un livello al suo interno
    alla ricerca del membro indicato, attraversando opzionalmente la
    catena di ereditarietà.
  \item \texttt{Call(query)} (\textit{Invocazione Transitiva}):
    Rappresenta l'esecuzione di un blocco di codice. Indica al
    valutatore che la sub-query passata in input identifica una
    funzione (o metodo), ma che l'espressione corrente necessita del
    suo \textbf{tipo di ritorno} per poter proseguire. È la primitiva
    fondamentale che abilita la concatenazione di metodi (es. in
      \texttt{a.f().g()}, il tipo di ritorno di \texttt{f()} diventa il
    \texttt{Find} di partenza per \texttt{g()}).
\end{itemize}

Attraverso questa algebra, un'espressione testuale complessa come
\texttt{i.f()} (con \texttt{i} di tipo \texttt{a.b.T}) viene espansa
nella query innestata:
$$ \texttt{Call(Extract(Extract(Extract(Find('a'), 'b'), 'T'), 'f'))} $$

Al termine dell'attraversamento, l'ambiente lessicale $\rho$ viene
completamente distrutto, lasciando l'IR popolato unicamente da
espressioni che non dipendono dalle variabili locali.

\section{Fase 2b: Esecuzione delle Query (Navigation Phase)}

La fase di navigazione ($R_{exec}$) si attua \textbf{solo dopo} aver
letto e tradotto tutto il codice dell'intero Workspace. Questa
posticipazione garantisce che l'ordine di dichiarazione delle classi
nei file sorgenti sia ininfluente.

Per valutare le query in modo deterministico e senza stati mutabili,
l'algoritmo si avvale di un \textbf{Contesto di Esecuzione
Immutabile} ($\Gamma$). Esso racchiude la conoscenza globale del
progetto e lo stato della traversata, ed è formalmente definito come la tupla:
$$ \Gamma = \langle \mathcal{P}_{curr}, \mathcal{I}_{stack},
\mathcal{GR}, \mathcal{PR} \rangle $$
dove:
\begin{itemize}
  \item $\mathcal{P}_{curr}$ è il prefisso gerarchico corrente (il
      percorso $\mathcal{QN}$ del componente in cui ci si trova, es.
    \texttt{root::core}).
  \item $\mathcal{I}_{stack}$ è la pila gerarchica degli import
    accessibili dallo scope attuale (i cosiddetti \textit{Jump Points}).
  \item $\mathcal{GR}$ è il \textbf{Global Registry}: un indice
    piatto creato a tempo zero, appena prima dell'esecuzione. Mappa
    ogni percorso assoluto ($\mathcal{QN}$) ai suoi metadati
    architetturali (es. tipi di ritorno e firme ereditate).
  \item $\mathcal{PR}$ è il \textbf{Primitive Registry}, il
    dizionario data-driven contenente le parole chiave dei tipi base
    del linguaggio in uso.
\end{itemize}

\subsection{Formalizzazione del Global Registry ($\mathcal{GR}$)}
Per permettere alla risoluzione di abbandonare l'oneroso approccio di
ricerca ad albero  in favore di un lookup piú leggero,  il
\textbf{Global Registry} è formalizzato come una mappa matematica
iniettiva tra i percorsi assoluti e i metadati strutturali essenziali.

$$ \mathcal{GR} : \mathcal{QN} \to E_{reg} $$
Dove $E_{reg}$ rappresenta il tipo somma delle entità architetturali
registrabili:
$$ E_{reg} = \text{Module} \mid \text{StructuredType} \mid
\text{Function}(\tau_{ret} \mid \mathcal{Q}_{ret}) \mid \text{ImplBlock} $$

L'intuizione architetturale fondamentale che distingue il
$\mathcal{GR}$ da un banale *hash set* di stringhe risiede nel
salvataggio del \textbf{tipo di ritorno} ($\tau_{ret}$ o la sua query
non risolta $\mathcal{Q}_{ret}$) per le entità funzionali.
Questa struttura dati permette al valutatore, incontrando l'operatore
di invocazione $\texttt{Call(F)}$, di non limitarsi a validare
l'esistenza della funzione $F$, ma di recuperarne istantaneamente il
tipo di ritorno interrogando $\mathcal{GR}(F)$. Ciò costituisce la
base algoritmica indispensabile per l'inferenza di tipo transitiva e
la validazione delle catene di metodi (es. $\texttt{a.f().g()}$).

\subsection{Lexical Scope Climbing e Punti di Salto}
Per valutare un operatore \texttt{Find}, il contesto di esecuzione
$\Gamma$ applica l'algoritmo di \textit{Lexical Scope Climbing}. Partendo
dal prefisso gerarchico corrente in cui la query è stata generata
(es. \texttt{root::core}), l'algoritmo accoda il nome cercato e
interroga il $\mathcal{GR}$.
Se la risorsa non viene trovata, il prefisso viene "tagliato" (pop)
risalendo di un livello (es. \texttt{root}), e il controllo viene
ripetuto. Contestualmente, a ogni livello risalito, l'algoritmo
ispeziona i \textbf{Jump Points} (le clausole di \texttt{import} o
\texttt{use}), verificando se il nome è un \textit{alias} che punta a
una risorsa esterna. Questo meccanismo riproduce fedelmente le regole
di \textit{Lexical Scoping} presenti nei compilatori reali.

\subsection{Method Resolution Order (MRO)}
L'operatore \texttt{Extract(P, m)} implementa nativamente la logica
di \textit{Dynamic Dispatch Statica}. Se il membro $m$ non esiste
direttamente nel perimetro del padre $P$, il valutatore estrae i
super-tipi (\texttt{super\_types}) di $P$ dal $\mathcal{GR}$ e lancia
una ricerca ricorsiva \textit{Depth-First, Left-to-Right} (Profondità
prima, da sinistra a destra) lungo l'albero genealogico.
Questo approccio risolve due problematiche architetturali:
\begin{enumerate}
  \item \textbf{Polimorfismo di Default}: Garantisce il corretto
    aggancio ai metodi definiti come default all'interno di Trait
    (Rust) o Interfacce (Java).
  \item \textbf{Diamond Problem}: Nei linguaggi che consentono
    l'ereditarietà multipla (C++, Python), se più padri espongono la
    stessa funzione, la discesa deterministica assicura l'estrazione
    della prima implementazione valida, salvaguardando l'integrità del grafo.
\end{enumerate}

\subsection{Inferenza Semantica e Decoppiamento Primitivo}
Durante l'esecuzione, il valutatore si avvale di ulteriori euristiche.
All'ingresso di un tipo strutturato, inietta le query relative ai
lessemi \texttt{self} e \texttt{this}, permettendo la risoluzione
immediata delle chiamate a metodi di classe. Inoltre, affronta il
problema dei falsi negativi operando un \textbf{Decoppiamento
Primitivo}: qualora un'intera query fallisca la sua valutazione nel
$\mathcal{GR}$, il sistema interroga il \textit{Primitive Registry}
(un dizionario data-driven caricato a compile-time). Se il nome base
coincide con una parola chiave del linguaggio (es. \texttt{int},
\texttt{String}), il riferimento viene declassato allo stato
\texttt{Prim}, evitando di contrassegnarlo come dipendenza
\texttt{Failed} o \texttt{External} e garantendo la pulizia semantica
del grafo architetturale finale.
